\documentclass{article}
\usepackage{amsmath,hyperref}
\title{A second fault changes the supplier of an unsettled duty}
\begin{document}
\maketitle
\paragraph{Pins and envelope.}
The native warehouse already committed an old keyed result, and no successor has yet installed its final checkpoint. The pinned SubT code supplies rolling byte accounting and registered delivery callbacks. Q and P have the same archived suffix. The source of the initial header is P at tick 6 regardless of the later mode. The later path change, notice truth, and authority rule are authored. At tick 7 exactly one sender can be prepared for the sole 173-byte suffix slot at tick 8. Mode \(Q\) keeps only Q's relevant path; mode \(P\) additionally loses Q's route to the suffix recipient and makes only the substitute P's path usable. Q can still serve the distinct diagnostic endpoint. Sender preparation is represented by a successful registered named request at tick 7. Unrelated warehouse work must run in both modes. Neither stock RWARE governance nor a measured SubT handoff protocol is claimed.

\paragraph{Common rule.}
Partition the joint histories by the actual callback available to the entry actor at tick 7. A candidate pair chooses the one requested sender and either tick 7 or tick 10 for the second diagnostic. The shared interpreter requires a licensed and physically supported sender, metered per-transmitter demand, one payment for the old duty, and both diagnostics by their own deadlines. It never reads the hidden physical mode. A discriminator selects one feasible candidate for the observed cell; when the cell has none, refuse the dependent action and preserve unrelated motion. The native effect also requires a delivered matching suffix, validated prefix and authorized install before it returns the keyed old result. Treating link replacement as renaming a sender would omit both old byte debits and its authority.

\paragraph{Exact scope.}
Q has already spent \(140+4\) bytes in the suffix window, so it can send \(173\) more but not also a second \(4\)-byte check before the boundary expires. P has \(114\) bytes from the header and additional debit \(d\); its suffix is admissible iff \(d+114+173\le320\), i.e.\ \(d\le33\), assuming its authority. Thus every admitted mode meets all duties with timely truthful notice, both licenses and \(d\le33\): prepare Q and defer the check in mode Q; prepare P and use the early check in mode P. A hidden selection between opposed physical modes cannot be implemented by a single prepared sender. A second preparation or redundant suffix can escape that information obstruction under a different authority or resource budget; it is not a legal move in this fixture. A late report does not move the already passed preparation cut. These are conditional guarantees, not an argument that the notification would naturally exist in every mission.

\paragraph{Testing and comparison.}
An exact direct two-budget calculation on the same observation cells and action catalogue agrees with the shared interpreter. The pinned native replay includes missing/late disclosure, one-byte currency of headroom, and asymmetric authority. Timely selected callbacks and actual suffix or refusal occur on one replay, alongside native unaffected movement. The exercise is a positive compound-fault development envelope and negative controls, not a full architecture coverage or comparative advantage claim.
\end{document}
